% !TeX program = xelatex
% ============================================================================
%  第 12 课　费马小定理的证明　—— 学生版讲义
%  与 02-课件/第12课-费马小定理/第12课-费马小定理.tex 同步
% ============================================================================

\xhlesson{12}{费马小定理的证明}{}

\begin{mubiao}
  \begin{itemize}
    \item 能说清为什么两边可以划掉相同的部分。
    \item 能从具体数字过渡到用字母写的一般形式。
    \item 能写出 \(a^{\,p-1}\equiv1\pmod p\)。
  \end{itemize}
\end{mubiao}

\section*{一、从一组数字开始}

\begin{dongshou}[取 \(p=7\)，\(a=3\)，把表填满]
  \begin{center}
  \begin{tabular}{@{}lcccccc@{}}
    \toprule
    原来的数 & \(1\) & \(2\) & \(3\) & \(4\) & \(5\) & \(6\) \\
    \midrule
    乘 \(3\) 之后 & \(3\) & \(6\) & \(9\) & \(12\) & \(15\) & \(18\) \\
    取余数 & & & & & & \\
    \bottomrule
  \end{tabular}
  \end{center}
\end{dongshou}

\begin{sikao}
  看最后一行：它是 \(1,2,3,4,5,6\) 的一次\underline{\hspace{5em}}？\\
  有没有哪个数字出现了两次？有没有哪个没出现？
  \liukong{2}
\end{sikao}

\begin{example}[把两行各自乘起来]
  \[
    1\times2\times3\times4\times5\times6=6!, \qquad
    3\times6\times2\times5\times1\times4=6! .
  \]
  因为下排只是上排的重排，两个乘积一模一样。
\end{example}

\section*{二、换一个角度乘}

\begin{dongshou}
  \[
    3\times6\times9\times12\times15\times18
    =\underline{\hspace{2.4em}}\times(1\times2\times3\times4\times5\times6)
    =\underline{\hspace{2.4em}}\times 6! .
  \]
  所以两边取模 \(7\) 得到：
  \[
    3^{6}\times6!\equiv6! \pmod 7 .
  \]
\end{dongshou}

\begin{sikao}[为什么能把 \(6!\) 划掉？]
  提示：模 \(7\) 里，什么样的数才有逆元？（想想第 3 课。）
  \liukong{2}
  因为 \(6!\) 里没有一个因数是 \(7\)，所以它\underline{\hspace{3em}}，两边可以同时划掉。\\[4pt]
  划完得到：\(3^{6}\equiv\underline{\hspace{3em}}\pmod 7\)。
\end{sikao}

\section*{三、换成字母再写一遍}

\begin{example}
  \begin{center}
  \begin{tabular}{@{}ll@{}}
    \toprule
    \textbf{具体}（\(p=7\)，\(a=3\)） & \textbf{一般} \\
    \midrule
    \(3,\ 6,\ 9,\ \dots,\ 18\) & \(a,\ 2a,\ 3a,\ \dots,\ (p-1)a\) \\
    余数是 \(1\text{—}6\) 的重排 & 余数是 \underline{\hspace{3.4em}} 的重排 \\
    \(3^{6}\times6!\equiv6!\) & \underline{\hspace{5.4em}} \\
    \(3^{6}\equiv1 \pmod 7\) & \underline{\hspace{5.4em}} \\
    \bottomrule
  \end{tabular}
  \end{center}
\end{example}

\begin{dingli}[费马小定理]
  设 \(p\) 是素数，\(a\) 是不被 \(p\) 整除的整数，那么
  \[
    a^{\,p-1}\equiv1 \pmod p .
  \]
\end{dingli}

\begin{sikao}
  整个证明用了两块砖，各是什么？
  \liukong{3}
\end{sikao}

\begin{xiaojie}
  不看上面的内容，凭记忆把证明的四个步骤写一遍：
  \liukong{5}
\end{xiaojie}

\noindent\textbf{下节课}：它有什么用——先去回答第 1 课黑板右边那个问题。